\documentclass[11pt,a4paper]{article}
\usepackage[utf8]{inputenc}
\usepackage[T1]{fontenc}
\usepackage[french]{babel}
\usepackage{amsmath, amssymb, amsthm}
\usepackage{geometry}
\geometry{hmargin=2.5cm, vmargin=2.5cm}
\usepackage{tikz}
\usepackage{listings}
\usepackage{xcolor}


\lstset{
    literate={é}{{\'e}}1 {ê}{{\^e}}1
             {è}{{\`e}}1
             {à}{{\`a}}1
             {ç}{{\c{c}}}1
             {î}{{\^i}}1
             {ù}{{\`u}}1
             {ô}{{\^o}}1
}
\lstset{
    language=Python,
    basicstyle=\ttfamily\small,
    keywordstyle=\color{blue},
    stringstyle=\color{red},
    commentstyle=\color{green!60!black},
    numbers=left,
    numberstyle=\tiny,
    stepnumber=1,
    numbersep=5pt,
    backgroundcolor=\color{gray!10},
    showspaces=false,
    showstringspaces=false,
    showtabs=false,
    frame=single,
    tabsize=4,
    captionpos=b,
    breaklines=true,
    breakatwhitespace=false,
    escapeinside={\%*}{*)}
}

\title{Endomorphismes symétriques, adjoint d'un opérateur et matrices orthogonales}
\author{Charles EDOU NZE}
\date{}

\begin{document}
\maketitle
\tableofcontents
\newpage




\section{Intuition et genèse du concept}

\begin{center}
\begin{tikzpicture}[scale=1.5]
    \draw[->, thick] (0,0) -- (2,0) node[right] {$x$};
    \draw[->, thick] (0,0) -- (0,2) node[above] {$y$};
    \draw[->, thick, red] (0,0) -- (1.5, 0.5) node[right] {$f(x)$};
    \draw[->, thick, blue] (0,0) -- (-0.5, 1.5) node[above] {$f^*(y)$};
    \draw[dashed] (1.5,0.5) -- (1.5,0);
    \draw[dashed] (-0.5,1.5) -- (0,1.5);
\end{tikzpicture}
\end{center}


L'algèbre linéaire a longtemps été considérée comme l'étude des transformations de l'espace abstrait. Mais lorsque l'on munit cet espace d'une géométrie stricte par le biais d'un produit scalaire (une "métrique"), de nouvelles questions vertigineuses émergent. Si je dispose d'une transformation (un endomorphisme) qui déforme l'espace, comment cette déformation interagit-elle avec la géométrie sous-jacente ? Existe-t-il des transformations qui "respectent" parfaitement l'angle et la distance, ou qui se déforment d'une manière si harmonieuse qu'elles conservent des directions privilégiées immuables ?

Historiquement, l'étude des endomorphismes symétriques trouve ses racines dans la mécanique céleste et l'étude des quadriques (coniques en dimension 2, ellipsoïdes en dimension 3). Augustin-Louis Cauchy, au début du XIXe siècle, cherchait à comprendre les axes principaux d'inertie des solides. Pourquoi un objet en rotation a-t-il toujours des axes naturels autour desquels il tourne sans "vaciller" ? La réponse mathématique est éblouissante de pureté : la matrice d'inertie est symétrique, et toute matrice symétrique, par un miracle algébrique, possède un ensemble de directions orthogonales qu'elle se contente d'étirer ou de compresser, sans jamais les faire pivoter. C'est l'essence même du célèbre théorème spectral.

Pour y parvenir, les mathématiciens ont dû inventer un concept miroir : l'adjoint d'un opérateur. L'adjoint, souvent noté $f^*($, est une construction intellectuelle fascinante. Imaginez que vous appliquez une transformation $f$ à un vecteur $x$, puis que vous mesurez sa "projection" (produit scalaire) sur un vecteur test $y$. L'opérateur adjoint $f^)$ est la machine qui vous permet d'obtenir \textit{exactement le même résultat de mesure} en laissant $x$ tranquille, mais en appliquant au préalable $f^*$ à $y$. C'est une dualité fondamentale. Un opérateur "symétrique" est simplement un opérateur qui est son propre miroir : l'effet qu'il a sur le monde est indiscernable de l'effet que le monde a sur lui. En Intelligence Artificielle, cette symétrie se retrouve partout, des matrices de covariance dans l'Analyse en Composantes Principales (PCA) aux matrices hessiennes de la fonction de coût, régissant la courbure de l'espace d'optimisation.

\section{2. Protocole d'Exégèse Conceptuelle}

\subsection{2.1. L'Adjoint d'un Opérateur (Endomorphisme)}

\textbf{A. Énoncé Symbolique Strict}

Soit $(E, \langle \cdot, \cdot \rangle)$ un espace euclidien (donc de dimension finie sur $\mathbb{R}$).
Pour tout endomorphisme $f \in \mathcal{L}(E)$, il existe un unique endomorphisme, noté $f^*$, appelé l'adjoint de $f$, tel que :
$$ \forall (x,y) \in E \times E, \quad \langle f(x), y \rangle = \langle x, f^*(y) \rangle $$

\textbf{B. Anatomie et Typage Chirurgical}
- $E$ : Un espace vectoriel sur le corps des réels $\mathbb{R}$, de dimension finie, muni d'un produit scalaire défini positif. L'hypothèse de dimension finie garantit l'existence de l'adjoint par le théorème de représentation de Riesz.
- $\langle \cdot, \cdot \rangle : E \times E \to \mathbb{R}$ : La forme bilinéaire symétrique définie positive équipant $E$.
- $f \in \mathcal{L}(E)$ : Un opérateur linéaire transformant les vecteurs de $E$ en d'autres vecteurs de $E$.
- $x, y \in E$ : Deux vecteurs quelconques sur lesquels s'opère la "mesure" d'interaction via le produit scalaire.
- L'égalité traduit un transfert de l'action de $f$ d'un côté du produit scalaire vers l'autre.

\textbf{C. Exemples de Validation}
- \textit{Exemple trivial :} Si $f = \text{Id}_E$ (l'identité), on a $\langle \text{Id}_E(x), y \rangle = \langle x, y \rangle = \langle x, \text{Id}_E(y) \rangle$. Donc l'adjoint de l'identité est elle-même : $\text{Id}_E^* = \text{Id}_E$.
- \textit{Exemple matriciel :} Si $E = \mathbb{R}^n$ muni du produit scalaire canonique $\langle X, Y \rangle = X^T Y$, et $f$ est représentée par la matrice $A$. Alors $\langle AX, Y \rangle = (AX)^T Y = X^T A^T Y = \langle X, A^T Y \rangle$. L'adjoint de $f$ est représenté par la matrice transposée $A^T$.

\textbf{D. Cas Pathologiques et Contre-exemples}
- \textit{Dimension infinie :} Si $E$ est un espace préhilbertien de dimension infinie, l'existence de l'adjoint n'est pas garantie pour un endomorphisme continu quelconque. Il faut se placer dans un espace de Hilbert (complet) pour s'appuyer sur le théorème de Riesz, et même là, l'adjoint d'opérateurs non bornés devient un sujet extrêmement subtil nécessitant la définition précise du domaine de définition $\mathcal{D}(f^*)$.

\subsection{2.2. Endomorphismes Symétriques}

\textbf{A. Énoncé Symbolique Strict}

Un endomorphisme $f \in \mathcal{L}(E)$ d'un espace euclidien est dit symétrique (ou autoadjoint) si :
$$ f = f^* \iff \forall (x,y) \in E^2, \quad \langle f(x), y \rangle = \langle x, f(y) \rangle $$
Dans une base orthonormée $\mathcal{B}$, la matrice $A = \text{Mat}_{\mathcal{B}}(f)$ est symétrique, c'est-à-dire $A^T = A$.

\textbf{B. Anatomie et Typage Chirurgical}
- $f = f^*$ : Égalité dans l'espace $\mathcal{L}(E)$. L'opérateur coïncide parfaitement avec son adjoint.
- Base orthonormée : Il est crucial de noter que la correspondance entre "endomorphisme symétrique" et "matrice symétrique" n'est vraie que si la base de projection est orthonormée. Dans une base quelconque, la transposée ne représente pas l'adjoint !

\textbf{C. Exemples de Validation}
- \textit{Exemple trivial :} L'endomorphisme nul et l'identité sont symétriques.
- \textit{Exemple géométrique :} La projection orthogonale $p$ sur un sous-espace $F$. On a $x = x_F + x_{F^\perp}$ et $y = y_F + y_{F^\perp}$.
$\langle p(x), y \rangle = \langle x_F, y_F + y_{F^\perp} \rangle = \langle x_F, y_F \rangle$.
De même $\langle x, p(y) \rangle = \langle x_F + x_{F^\perp}, y_F \rangle = \langle x_F, y_F \rangle$. Donc $p = p^*$.

\textbf{D. Cas Pathologiques et Contre-exemples}
- Une projection oblique (non orthogonale) n'est jamais symétrique. L'asymétrie de la projection (qui "pousse" l'espace de biais) détruit la propriété miroir vis-à-vis du produit scalaire canonique.

\subsection{2.3. Matrices et Transformations Orthogonales (Isométries)}

\textbf{A. Énoncé Symbolique Strict}

Un endomorphisme $u \in \mathcal{L}(E)$ est orthogonal (ou une isométrie vectorielle) si :
$$ \forall x \in E, \quad \|u(x)\| = \|x\| $$
Une matrice $O \in \mathcal{M}_n(\mathbb{R})$ est orthogonale si :
$$ O^T O = O O^T = I_n $$

\textbf{B. Anatomie et Typage Chirurgical}
- $\| \cdot \|$ : La norme induite par le produit scalaire. La conservation de la norme implique (par identité de polarisation) la conservation du produit scalaire : $\langle u(x), u(y) \rangle = \langle x, y \rangle$.
- $O^T O = I_n$ : Les colonnes (et les lignes) de $O$ forment une base orthonormée de $\mathbb{R}^n$.

\textbf{C. Exemples de Validation}
- \textit{Rotations et Réflexions :} En dimension 2, la matrice $R_\theta = \begin{pmatrix} \cos\theta & -\sin\theta \\ \sin\theta & \cos\theta \end{pmatrix}$ satisfait $R_\theta^T R_\theta = I_2$. C'est une isométrie qui conserve les distances et les angles.

\textbf{D. Cas Pathologiques et Contre-exemples}
- Matrice orthogonale vs matrice symétrique : Attention à ne pas confondre. $R_{\pi/2}$ est orthogonale mais pas symétrique. Inversement, $2I_n$ est symétrique mais pas orthogonale (elle étire les vecteurs). Les matrices à la fois symétriques et orthogonales sont les symétries orthogonales ($s^2 = Id$ et $s^* = s$).

\section{3. Zéro Ellipse : Preuves Exhaustives}

\subsection{3.1. Théorème d'existence et unicité de l'adjoint}

\textbf{Énoncé :} Pour tout $f \in \mathcal{L}(E)$, il existe un unique $f^*($ tel que $\forall x, y, \langle f(x), y \rangle = \langle x, f^)(y) \rangle$.

\textbf{Démonstration :}
1. Fixons un vecteur $y \in E$.
2. Considérons l'application $\varphi_y : E \to \mathbb{R}$ définie par $\varphi_y(x) = \langle f(x), y \rangle$.
3. Montrons que $\varphi_y$ est une forme linéaire sur $E$.
   Soient $x_1, x_2 \in E$ et $\lambda \in \mathbb{R}$.
   $\varphi_y(\lambda x_1 + x_2) = \langle f(\lambda x_1 + x_2), y \rangle$.
   Par linéarité de $f$, $f(\lambda x_1 + x_2) = \lambda f(x_1) + f(x_2)$.
   Par bilinéarité (linéarité à gauche) du produit scalaire,
   $\langle \lambda f(x_1) + f(x_2), y \rangle = \lambda \langle f(x_1), y \rangle + \langle f(x_2), y \rangle = \lambda \varphi_y(x_1) + \varphi_y(x_2)$.
   Donc $\varphi_y \in E^*$.
4. D'après le théorème de représentation de Riesz (qui stipule que l'application $z \mapsto \langle \cdot, z \rangle$ est un isomorphisme de $E$ sur $E^*$), il existe un unique vecteur $z_y \in E$ tel que pour tout $x \in E$, $\varphi_y(x) = \langle x, z_y \rangle$.
5. On définit alors l'application $f^*( : E \to E$ par $f^)(y) = z_y$. Par construction, $\langle f(x), y \rangle = \langle x, f^*(y) \rangle$.
6. Montrons que $f^*$ est linéaire. Soient $y_1, y_2 \in E$ et $\alpha \in \mathbb{R}$. Pour tout $x \in E$, évaluons le produit scalaire :
   $\langle x, f^*(\alpha y_1 + y_2) \rangle = \langle f(x), \alpha y_1 + y_2 \rangle$.
   Par linéarité à droite du produit scalaire (rappelons que sur $\mathbb{R}$, bilinéaire implique linéaire des deux côtés),
   $\langle f(x), \alpha y_1 + y_2 \rangle = \alpha \langle f(x), y_1 \rangle + \langle f(x), y_2 \rangle$.
   En appliquant la définition de $f^*$ aux deux termes :
   $\alpha \langle f(x), y_1 \rangle + \langle f(x), y_2 \rangle = \alpha \langle x, f^*((y_1) \rangle + \langle x, f^)(y_2) \rangle = \langle x, \alpha f^*((y_1) + f^)(y_2) \rangle$.
   Ainsi, $\forall x \in E, \langle x, f^*((\alpha y_1 + y_2) - (\alpha f^)(y_1) + f^*(y_2)) \rangle = 0$.
   Le seul vecteur orthogonal à tout l'espace étant le vecteur nul, on en déduit formellement $f^*((\alpha y_1 + y_2) = \alpha f^)(y_1) + f^*(y_2)$.
   $f^*$ est donc bien un endomorphisme.

\subsection{3.2. Propriété d'orthogonalité des sous-espaces propres (Symétriques)}

\textbf{Énoncé :} Si $f$ est un endomorphisme symétrique, alors les sous-espaces propres associés à des valeurs propres distinctes sont deux à deux orthogonaux.

\textbf{Démonstration :}
1. Soient $\lambda, \mu \in \mathbb{R}$ deux valeurs propres distinctes de $f$ ($\lambda \neq \mu$).
2. Soient $x \in E_\lambda$ et $y \in E_\mu$ deux vecteurs propres associés, c'est-à-dire $f(x) = \lambda x$ et $f(y) = \mu y$.
3. Calculons la quantité $\langle f(x), y \rangle$.
   D'une part, en remplaçant $f(x)$ :
   $\langle f(x), y \rangle = \langle \lambda x, y \rangle$.
   Par linéarité à gauche du produit scalaire :
   $\langle \lambda x, y \rangle = \lambda \langle x, y \rangle$.
4. D'autre part, en utilisant la symétrie de l'endomorphisme $f$ ($f = f^*$) :
   $\langle f(x), y \rangle = \langle x, f^*(y) \rangle = \langle x, f(y) \rangle$.
   En remplaçant $f(y)$ par $\mu y$ :
   $\langle x, \mu y \rangle = \mu \langle x, y \rangle$ (par symétrie et linéarité à gauche du produit scalaire).
5. En égalisant les deux résultats obtenus :
   $\lambda \langle x, y \rangle = \mu \langle x, y \rangle$.
   En soustrayant :
   $(\lambda - \mu) \langle x, y \rangle = 0$.
6. Puisque par hypothèse les valeurs propres sont distinctes ($\lambda - \mu \neq 0$), l'intégrité du corps des réels impose nécessairement que le second facteur soit nul :
   $\langle x, y \rangle = 0$.
7. Les vecteurs $x$ et $y$ sont donc orthogonaux. Par extension, tous les vecteurs de $E_\lambda$ sont orthogonaux à tous les vecteurs de $E_\mu$, ce qui prouve l'orthogonalité des sous-espaces.

\subsection{3.3. Théorème Spectral (Diagonalisabilité des Endomorphismes Symétriques)}

\textbf{Énoncé :} Tout endomorphisme symétrique d'un espace euclidien $E$ possède une base orthonormée de vecteurs propres. En d'autres termes, toute matrice réelle symétrique est diagonalisable en base orthonormée : $A = P D P^T$ avec $P$ orthogonale.

\textit{(La preuve complète requiert l'existence d'au moins une valeur propre réelle, généralement démontrée par l'analyse sur les polynômes caractéristiques ou la compacité de la sphère unité, détaillée dans le jalon suivant. Nous admettons ici l'existence d'au moins une valeur propre $\lambda_1$ et d'un vecteur propre unitaire $e_1$).}

\textbf{Démonstration par récurrence sur la dimension $n$ :}
1. \textbf{Initialisation ($n=1$) :} L'espace est de dimension 1. Tout endomorphisme est représenté par un scalaire, et toute base unitaire $(e_1)$ est orthogonale. La propriété est triviale.
2. \textbf{Hérédité :} Supposons le théorème vrai pour tout espace de dimension $n-1$. Soit $E$ un espace de dimension $n$, et $f$ un endomorphisme symétrique sur $E$.
3. On admet que $f$ possède au moins un vecteur propre $e_1$ associé à la valeur propre $\lambda_1$. On normalise ce vecteur tel que $\|e_1\| = 1$.
4. Posons $H = (\text{Vect}(e_1))^\perp$, l'hyperplan orthogonal à $e_1$. $\dim(H) = n - 1$.
5. Montrons que $H$ est stable par $f$.
   Soit $x \in H$. On veut montrer que $f(x) \in H$, c'est-à-dire $\langle f(x), e_1 \rangle = 0$.
   Calculons : $\langle f(x), e_1 \rangle$.
   Par symétrie de $f$ : $\langle f(x), e_1 \rangle = \langle x, f(e_1) \rangle$.
   Puisque $e_1$ est vecteur propre de $f$ : $f(e_1) = \lambda_1 e_1$.
   Donc $\langle x, \lambda_1 e_1 \rangle = \lambda_1 \langle x, e_1 \rangle$.
   Or, $x \in H = (\text{Vect}(e_1))^\perp$, ce qui signifie par définition que $\langle x, e_1 \rangle = 0$.
   Ainsi, $\langle f(x), e_1 \rangle = \lambda_1 \times 0 = 0$.
   Le vecteur $f(x)$ est bien orthogonal à $e_1$, donc $f(x) \in H$. $H$ est stable par $f$.
6. Considérons l'endomorphisme restreint $f_H : H \to H$ défini par $f_H(x) = f(x)$.
   $f_H$ est clairement linéaire et symétrique (car $f$ l'est sur tout $E$, donc sur $H$).
7. Par hypothèse de récurrence, puisque $\dim(H) = n-1$, l'endomorphisme $f_H$ possède une base orthonormée de vecteurs propres dans $H$, notons-la $(e_2, e_3, \dots, e_n)$.
8. En rassemblant ces vecteurs avec $e_1$, la famille $(e_1, e_2, \dots, e_n)$ est constituée de vecteurs propres de $f$.
   Par construction, les $e_2, \dots, e_n$ appartiennent à $H = \{e_1\}^\perp$, donc ils sont tous orthogonaux à $e_1$.
   De plus, $(e_2, \dots, e_n)$ est une famille orthonormée.
   Donc la famille totale $(e_1, e_2, \dots, e_n)$ est orthonormée.
9. C'est une famille orthonormée de $n$ vecteurs dans un espace de dimension $n$, c'est donc une base orthonormée.
10. La récurrence est établie, et le théorème spectral est rigoureusement prouvé.

\section{Exercices d'application et de concours}

\subsection*{Exercice 01 : Adjoint d'une composée}


\textbf{Énoncé :}


Démontrer que pour tout $f, g \in \mathcal{L}(E)$, on a $(f \circ g)^* = g^* \circ f^*$.

\textbf{Démonstration :}


Soient $x, y \in E$. Par définition de l'adjoint de $f \circ g$, on a :
$$ \langle (f \circ g)(x), y \rangle = \langle x, (f \circ g)^*(y) \rangle $$
D'autre part, en développant la composition :
$$ \langle (f \circ g)(x), y \rangle = \langle f(g(x)), y \rangle $$
En appliquant la définition de l'adjoint de $f$ à la paire $(g(x), y)$ :
$$ \langle f(g(x)), y \rangle = \langle g(x), f^*(y) \rangle $$
Ensuite, en appliquant la définition de l'adjoint de $g$ à la paire $(x, f^*(y))$ :
$$ \langle g(x), f^*(y) \rangle = \langle x, g^*(f^*(y)) \rangle = \langle x, (g^* \circ f^*)(y) \rangle $$
Ainsi, pour tout $x, y \in E$, on obtient :
$$ \langle x, (f \circ g)^*(y) \rangle = \langle x, (g^* \circ f^*)(y) \rangle $$
Ceci implique que $(f \circ g)^*(y) = (g^* \circ f^*)(y)$ pour tout $y \in E$.
Par conséquent, $(f \circ g)^* = g^* \circ f^*$. $\blacksquare$

\subsection*{Exercice 02 : Adjoint et inverse}


\textbf{Énoncé :}


Montrer que si un endomorphisme $f \in \mathcal{L}(E)$ est inversible, alors $f^*$ est inversible et $(f^*)^{-1} = (f^{-1})^*$.

\textbf{Démonstration :}


Soit $f \in \mathcal{L}(E)$ un endomorphisme inversible.
Par définition de l'inverse, nous avons $f \circ f^{-1} = f^{-1} \circ f = \operatorname{Id}_E$.
Nous savons que l'adjoint de l'identité est l'identité elle-même, car :
$$ \langle \operatorname{Id}_E(x), y \rangle = \langle x, y \rangle = \langle x, \operatorname{Id}_E(y) \rangle $$
Donc $\operatorname{Id}_E^* = \operatorname{Id}_E$.
Prenons l'adjoint de l'égalité $f \circ f^{-1} = \operatorname{Id}_E$. En utilisant la propriété de l'exercice 1, $(u \circ v)^* = v^* \circ u^*$, nous obtenons :
$$ (f \circ f^{-1})^* = \operatorname{Id}_E^* \implies (f^{-1})^* \circ f^* = \operatorname{Id}_E $$
De même, en prenant l'adjoint de $f^{-1} \circ f = \operatorname{Id}_E$ :
$$ (f^{-1} \circ f)^* = \operatorname{Id}_E^* \implies f^* \circ (f^{-1})^* = \operatorname{Id}_E $$
Les relations $(f^{-1})^* \circ f^* = \operatorname{Id}_E$ et $f^* \circ (f^{-1})^* = \operatorname{Id}_E$ démontrent exactement que l'endomorphisme $f^*$ est inversible, et que son inverse est donné par $(f^{-1})^*$. $\blacksquare$

\subsection*{Exercice 03 : Somme de projecteurs orthogonaux}


\textbf{Énoncé :}


Soient $P$ et $Q$ deux projecteurs orthogonaux sur $E$. À quelle condition nécessaire et suffisante $P+Q$ est-il un projecteur orthogonal ?

\textbf{Démonstration :}


Par définition, un projecteur orthogonal $R$ vérifie $R^2 = R$ (idempotence) et $R^* = R$ (symétrie).
Puisque $P$ et $Q$ sont des projecteurs orthogonaux, on a $P^2=P$, $P^*=P$, et $Q^2=Q$, $Q^*=Q$.
L'opérateur $P+Q$ est toujours symétrique, car :
$$ (P+Q)^* = P^* + Q^* = P + Q $$
La question se réduit donc à savoir quand $P+Q$ est un projecteur, c'est-à-dire quand $(P+Q)^2 = P+Q$.
Développons $(P+Q)^2$ :
$$ (P+Q)^2 = (P+Q) \circ (P+Q) = P^2 + P \circ Q + Q \circ P + Q^2 $$
En utilisant $P^2=P$ et $Q^2=Q$, cela donne :
$$ (P+Q)^2 = P + P \circ Q + Q \circ P + Q = (P+Q) + P \circ Q + Q \circ P $$
Pour que $(P+Q)^2 = P+Q$, il faut et il suffit que :
$$ P \circ Q + Q \circ P = 0 \quad \text{(Équation 1)} $$
Multiplions l'Équation 1 à droite par $P$ :
$$ P \circ Q \circ P + Q \circ P^2 = 0 \implies P \circ Q \circ P + Q \circ P = 0 \quad \text{(car } P^2=P) $$
Multiplions l'Équation 1 à gauche par $P$ :
$$ P^2 \circ Q + P \circ Q \circ P = 0 \implies P \circ Q + P \circ Q \circ P = 0 $$
En soustrayant ces deux nouvelles relations :
$$ (P \circ Q + P \circ Q \circ P) - (P \circ Q \circ P + Q \circ P) = 0 \implies P \circ Q - Q \circ P = 0 \implies P \circ Q = Q \circ P $$
Remplaçons $Q \circ P$ par $P \circ Q$ dans l'Équation 1 :
$$ 2(P \circ Q) = 0 \implies P \circ Q = 0 $$
Par symétrie, $Q \circ P = 0$.
Ainsi, la condition nécessaire et suffisante est que $P \circ Q = Q \circ P = 0$, ce qui signifie géométriquement que les images de $P$ et $Q$ sont orthogonales. $\blacksquare$

\subsection*{Exercice 04 : Endomorphisme antisymétrique}


\textbf{Énoncé :}


Un endomorphisme $f \in \mathcal{L}(E)$ est dit antisymétrique si $f^* = -f$. Montrer que si $f$ est antisymétrique, alors pour tout vecteur $x \in E$, on a $\langle f(x), x \rangle = 0$.

\textbf{Démonstration :}


Soit $f \in \mathcal{L}(E)$ un endomorphisme antisymétrique. Par définition, nous avons la relation $f^* = -f$.
Considérons le produit scalaire $\langle f(x), x \rangle$ pour un vecteur arbitraire $x \in E$.
D'après la définition de l'endomorphisme adjoint, nous pouvons transférer l'opérateur $f$ sur le second argument du produit scalaire :
$$ \langle f(x), x \rangle = \langle x, f^*(x) \rangle $$
En utilisant l'hypothèse d'antisymétrie $f^* = -f$, nous remplaçons $f^*(x)$ par $-f(x)$ :
$$ \langle x, f^*(x) \rangle = \langle x, -f(x) \rangle $$
Par bilinéarité du produit scalaire (spécifiquement la linéarité par rapport à la seconde variable), nous pouvons sortir le scalaire $-1$ :
$$ \langle x, -f(x) \rangle = -\langle x, f(x) \rangle $$
Par ailleurs, le produit scalaire sur un espace euclidien réel est symétrique, ce qui signifie que $\langle x, y \rangle = \langle y, x \rangle$ pour tout $x, y$. En appliquant cela, nous obtenons :
$$ -\langle x, f(x) \rangle = -\langle f(x), x \rangle $$
Nous avons ainsi établi l'égalité :
$$ \langle f(x), x \rangle = -\langle f(x), x \rangle $$
En additionnant $\langle f(x), x \rangle$ des deux côtés de l'équation, il vient :
$$ 2\langle f(x), x \rangle = 0 $$
Puisque le scalaire $2$ est non nul, il en résulte nécessairement que :
$$ \langle f(x), x \rangle = 0 $$
Cette propriété démontre que pour un endomorphisme antisymétrique, le vecteur image $f(x)$ est toujours orthogonal au vecteur source $x$. $\blacksquare$

\subsection*{Exercice 05 : Polynôme d'un endomorphisme symétrique}


\textbf{Énoncé :}


Soit $f$ un endomorphisme symétrique. Montrer que pour tout polynôme $P \in \mathbb{R}[X]$, l'endomorphisme $P(f)$ est également symétrique.

\textbf{Démonstration :}


Soit $f \in \mathcal{L}(E)$ un endomorphisme symétrique. Par définition, $f^* = f$.
Soit $P \in \mathbb{R}[X]$ un polynôme quelconque. Il s'écrit sous la forme $P(X) = \sum_{k=0}^d a_k X^k$, où $a_k \in \mathbb{R}$ sont ses coefficients et $d$ son degré.
L'endomorphisme $P(f)$ est défini par :
$$ P(f) = \sum_{k=0}^d a_k f^k $$
où $f^k = f \circ f \circ \dots \circ f$ ($k$ fois) et $f^0 = \operatorname{Id}_E$.
Calculons l'adjoint de $P(f)$. L'opération d'adjonction est linéaire, donc :
$$ (P(f))^* = \left( \sum_{k=0}^d a_k f^k \right)^* = \sum_{k=0}^d a_k (f^k)^* $$
Calculons l'adjoint de $f^k$. Par récurrence, en utilisant la propriété $(u \circ v)^* = v^* \circ u^*$ :
- Pour $k=0$, $(\operatorname{Id}_E)^* = \operatorname{Id}_E = f^0$.
- Pour $k=1$, $f^* = f = f^1$.
- Pour un $k$ quelconque, $(f^k)^* = (f \circ f^{k-1})^* = (f^{k-1})^* \circ f^* = (f^{k-1}) \circ f = f^k$.
Ainsi, pour tout $k \in \mathbb{N}$, on a $(f^k)^* = f^k$.
En réinjectant ce résultat dans l'expression de l'adjoint de $P(f)$ :
$$ (P(f))^* = \sum_{k=0}^d a_k f^k = P(f) $$
L'endomorphisme $P(f)$ est donc égal à son adjoint. Il est bien symétrique. $\blacksquare$

\subsection*{Exercice 06 : Valeurs propres réelles}


\textbf{Énoncé :}


Soit $A \in \mathcal{M}_n(\mathbb{R})$ une matrice symétrique. Montrer que ses valeurs propres complexes sont nécessairement réelles.

\textbf{Démonstration :}


Soit $A \in \mathcal{M}_n(\mathbb{R})$ symétrique. On la considère comme une matrice de $\mathcal{M}_n(\mathbb{C})$.
Puisque $A$ est réelle, $\bar{A} = A$. Puisque $A$ est symétrique, $A^T = A$.
Soit $\lambda \in \mathbb{C}$ une valeur propre de $A$, et $X \in \mathbb{C}^n$ un vecteur propre associé non nul.
On a $A X = \lambda X$.
Prenons le transconjugué (adjoint complexe $X^* = \bar{X}^T$) de l'équation $AX = \lambda X$.
Multiplions à gauche par $X^*$ :
$$ X^* (A X) = X^* (\lambda X) = \lambda X^* X $$
Le terme $X^* X = \sum_{i=1}^n |x_i|^2 > 0$ car $X \neq 0$.
Par ailleurs, prenons l'équation conjuguée de $AX = \lambda X$ :
$$ \overline{A X} = \overline{\lambda X} \implies \bar{A} \bar{X} = \bar{\lambda} \bar{X} $$
Comme $A$ est réelle, $A \bar{X} = \bar{\lambda} \bar{X}$.
Transposons cette équation :
$$ (A \bar{X})^T = (\bar{\lambda} \bar{X})^T \implies \bar{X}^T A^T = \bar{\lambda} \bar{X}^T $$
Comme $A$ est symétrique, $A^T = A$. Donc :
$$ \bar{X}^T A = \bar{\lambda} \bar{X}^T \implies X^* A = \bar{\lambda} X^* $$
Multiplions cette équation à droite par $X$ :
$$ (X^* A) X = (\bar{\lambda} X^*) X = \bar{\lambda} X^* X $$
Nous avons ainsi deux expressions pour le scalaire $X^* A X$ :
1. $X^* A X = \lambda X^* X$
2. $X^* A X = \bar{\lambda} X^* X$
En égalant les deux expressions, nous obtenons :
$$ \lambda X^* X = \bar{\lambda} X^* X $$
$$ (\lambda - \bar{\lambda}) X^* X = 0 $$
Puisque $X^* X \neq 0$, on en déduit que $\lambda - \bar{\lambda} = 0$, soit $\lambda = \bar{\lambda}$.
La valeur propre $\lambda$ est donc un nombre réel. $\blacksquare$

\subsection*{Exercice 07 : Diagonalisation simultanée}


\textbf{Énoncé :}


Soient $f$ et $g$ deux endomorphismes symétriques de $E$ qui commutent ($f \circ g = g \circ f$). Montrer qu'ils sont co-diagonalisables dans une même base orthonormée.

\textbf{Démonstration :}


Soient $f$ et $g$ deux endomorphismes symétriques qui commutent. D'après le théorème spectral, $f$ est diagonalisable. Ses sous-espaces propres $E_\lambda(f)$ sont mutuellement orthogonaux et leur somme directe est $E$.
Soit $\lambda$ une valeur propre de $f$. Montrons que $E_\lambda(f)$ est stable par $g$.
Soit $x \in E_\lambda(f)$. Alors $f(x) = \lambda x$.
Calculons $f(g(x))$. Comme $f$ et $g$ commutent, $f(g(x)) = g(f(x))$.
Puisque $f(x) = \lambda x$, on a $g(f(x)) = g(\lambda x) = \lambda g(x)$.
Donc $f(g(x)) = \lambda g(x)$, ce qui prouve que $g(x) \in E_\lambda(f)$.
Ainsi, la restriction $g_\lambda$ de $g$ à $E_\lambda(f)$ est un endomorphisme symétrique de l'espace euclidien $E_\lambda(f)$.
En appliquant le théorème spectral à $g_\lambda$ sur $E_\lambda(f)$, on trouve une base orthonormée $B_\lambda$ de $E_\lambda(f)$ formée de vecteurs propres de $g$.
Ces vecteurs sont par construction des vecteurs propres de $f$ (puisqu'ils sont dans $E_\lambda(f)$) et des vecteurs propres de $g$.
En réunissant les bases $B_\lambda$ pour toutes les valeurs propres $\lambda$ de $f$, on obtient une base $B$ de $E$.
Comme les sous-espaces propres $E_\lambda(f)$ sont orthogonaux deux à deux, et que chaque $B_\lambda$ est orthonormée, la réunion $B$ est une base orthonormée de $E$.
Dans cette base $B$, les matrices de $f$ et de $g$ sont simultanément diagonales. $\blacksquare$

\subsection*{Exercice 08 : Racine carrée symétrique définie positive}


\textbf{Énoncé :}


Soit $S$ un endomorphisme symétrique défini positif (toutes ses valeurs propres sont strictement positives). Montrer qu'il existe un unique endomorphisme symétrique défini positif $R$ tel que $R^2 = S$.

\textbf{Démonstration :}


\textbf{Existence :}
D'après le théorème spectral, il existe une base orthonormée $(e_1, \dots, e_n)$ de vecteurs propres de $S$ avec les valeurs propres associées $\lambda_1, \dots, \lambda_n$. Comme $S$ est défini positif, $\lambda_i > 0$ pour tout $i$.
On définit l'endomorphisme $R$ par $R(e_i) = \sqrt{\lambda_i} e_i$ pour tout $i$.
Par construction, $R$ est diagonalisable dans la base orthonormée $(e_1, \dots, e_n)$ avec des valeurs propres $\sqrt{\lambda_i} > 0$. Donc $R$ est symétrique et défini positif.
De plus, $R^2(e_i) = R(\sqrt{\lambda_i} e_i) = \sqrt{\lambda_i} R(e_i) = (\sqrt{\lambda_i})^2 e_i = \lambda_i e_i = S(e_i)$.
Les endomorphismes $R^2$ et $S$ coïncident sur une base, ils sont donc égaux : $R^2 = S$.

\textbf{Unicité :}
Soit $R'$ un autre endomorphisme symétrique défini positif tel que $(R')^2 = S$.
Comme $R'$ et $S$ commutent ($R' \circ S = R' \circ (R')^2 = (R')^3 = (R')^2 \circ R' = S \circ R'$), $R'$ conserve les sous-espaces propres de $S$.
Soit $E_\lambda$ le sous-espace propre de $S$ pour la valeur propre $\lambda$.
La restriction $r'$ de $R'$ à $E_\lambda$ est symétrique définie positive et vérifie $(r')^2 = \lambda \operatorname{Id}$.
Les valeurs propres de $r'$ doivent être racines de $X^2 - \lambda$. Comme $r'$ est défini positif, sa seule valeur propre possible est $\sqrt{\lambda}$.
Puisque $r'$ est symétrique avec une unique valeur propre $\sqrt{\lambda}$, $r'$ est obligatoirement l'homothétie de rapport $\sqrt{\lambda}$ sur $E_\lambda$.
Ainsi, $R'$ agit exactement comme $R$ sur chaque sous-espace propre de $S$.
Comme $E$ est somme directe orthogonale de ces sous-espaces, $R' = R$. L'unicité est prouvée. $\blacksquare$

\subsection*{Exercice 09 : Matrices orthogonales et isométrie}


\textbf{Énoncé :}


Montrer qu'une matrice carrée $A \in \mathcal{M}_n(\mathbb{R})$ est orthogonale si et seulement si l'endomorphisme canoniquement associé conserve la norme (c'est une isométrie).

\textbf{Démonstration :}


Soit $f$ l'endomorphisme associé à la matrice $A$ dans la base canonique, qui est orthonormée pour le produit scalaire standard $\langle X, Y \rangle = X^T Y$.
Dire que $A$ est orthogonale signifie que $A^T A = I_n$.
Dire que $f$ conserve la norme signifie que pour tout $X \in \mathbb{R}^n$, $\|AX\| = \|X\|$.

\textbf{Sens direct :} Supposons $A^T A = I_n$.
Pour tout $X \in \mathbb{R}^n$, la norme au carré est :
$$ \|AX\|^2 = \langle AX, AX \rangle = (AX)^T (AX) = X^T A^T A X $$
Puisque $A^T A = I_n$, cela devient :
$$ X^T I_n X = X^T X = \|X\|^2 $$
Ainsi, $\|AX\| = \|X\|$, l'endomorphisme conserve la norme.

\textbf{Sens réciproque :} Supposons que pour tout $X \in \mathbb{R}^n$, $\|AX\| = \|X\|$.
Alors pour tout $X$, $X^T A^T A X = X^T X$, ce qui s'écrit $X^T (A^T A - I_n) X = 0$.
Posons $M = A^T A - I_n$. La matrice $M$ est symétrique car $M^T = (A^T A - I_n)^T = A^T A - I_n = M$.
La forme quadratique associée à $M$ est nulle : pour tout $X$, $X^T M X = 0$.
Pour une matrice symétrique, si la forme quadratique est identiquement nulle, alors la matrice est nulle (cela se prouve par polarisation : $2 X^T M Y = (X+Y)^T M (X+Y) - X^T M X - Y^T M Y = 0 - 0 - 0 = 0$).
Donc $M = 0$, c'est-à-dire $A^T A = I_n$. La matrice $A$ est orthogonale. $\blacksquare$

\subsection*{Exercice 10 : Décomposition polaire}

Difficulté : $\star$$\star$$\star$$\star$$\star$

\textbf{Énoncé :}


Soit $A \in \mathcal{GL}_n(\mathbb{R})$ une matrice inversible. Démontrer le théorème de décomposition polaire : il existe un unique couple $(O, S)$ avec $O$ matrice orthogonale et $S$ matrice symétrique définie positive tel que $A = O S$.

\textbf{Démonstration :}


\textbf{Existence :}
Considérons la matrice $M = A^T A$. $M$ est symétrique, car $M^T = (A^T A)^T = A^T (A^T)^T = A^T A = M$.
De plus, $M$ est définie positive. Pour tout vecteur non nul $X \in \mathbb{R}^n$ :
$$ X^T M X = X^T (A^T A) X = (AX)^T (AX) = \|AX\|^2 $$
Puisque $A$ est inversible et $X \neq 0$, $AX \neq 0$, donc $\|AX\|^2 > 0$. Ainsi, $M$ est symétrique définie positive.
D'après l'exercice 8, il existe une unique matrice $S$ symétrique définie positive telle que $S^2 = M = A^T A$.
Puisque $S$ est définie positive, elle est inversible. Définissons $O = A S^{-1}$.
Il reste à vérifier que $O$ est une matrice orthogonale.
$$ O^T O = (A S^{-1})^T (A S^{-1}) = (S^{-1})^T A^T A S^{-1} $$
Puisque $S$ est symétrique, $S^{-1}$ est aussi symétrique, donc $(S^{-1})^T = S^{-1}$.
De plus, par construction, $A^T A = S^2$.
$$ O^T O = S^{-1} S^2 S^{-1} = S^{-1} S S S^{-1} = I_n I_n = I_n $$
La matrice $O$ est donc bien orthogonale. Nous avons $A = O S$ avec les propriétés requises.

\textbf{Unicité :}
Supposons que $A = O_1 S_1 = O_2 S_2$ soient deux décompositions polaires.
Alors $A^T A = (O_1 S_1)^T (O_1 S_1) = S_1^T O_1^T O_1 S_1 = S_1 I_n S_1 = S_1^2$.
De même, $A^T A = S_2^2$.
Donc $S_1^2 = S_2^2 = A^T A$.
Or, d'après l'exercice 8, la racine carrée symétrique définie positive d'une matrice symétrique définie positive est unique.
Donc $S_1 = S_2$.
En multipliant par $S_1^{-1}$ à droite, on obtient $O_1 = O_2$.
L'unicité est prouvée. $\blacksquare$

\section{Travaux pratiques et simulations algorithmiques}
\subsection*{TP 01 : Implémentation de la transposition et vérification de la symétrie}

\textbf{Objectif :}

Implémenter la transposition matricielle à partir de zéro et concevoir une fonction validant rigoureusement si un endomorphisme représenté par sa matrice est symétrique (i.e. $A = A^T$).



\begin{lstlisting}
def transpose(A):
    rows = len(A)
    cols = len(A[0])
    return [[A[i][j] for i in range(rows)] for j in range(cols)]

def is_symmetric(A):
    rows = len(A)
    cols = len(A[0])
    if rows != cols:
        return False
    for i in range(rows):
        for j in range(cols):
            # Utilisation d'une tolérance pour contrer les erreurs flottantes
            if abs(A[i][j] - A[j][i]) > 1e-9:
                return False
    return True

# Validation
A_sym = [[1.0, 2.0, 3.0],
         [2.0, 4.0, -1.0],
         [3.0, -1.0, 5.0]]

A_asym = [[1.0, 2.0, 3.0],
          [0.0, 4.0, -1.0],
          [3.0, -1.0, 5.0]]

assert is_symmetric(A_sym) == True, "La matrice A_sym doit être symétrique."
assert is_symmetric(A_asym) == False, "La matrice A_asym ne doit pas être symétrique."
assert transpose(A_sym) == A_sym, "La transposée d'une matrice symétrique est elle-même."

print("TP 01 validé : L'opérateur de transposition et la vérification de symétrie sont opérationnels.")
\end{lstlisting}

\subsection*{TP 02 : Vérification de la propriété de l'adjoint sur le produit scalaire}

\textbf{Objectif :}

Valider algorithmiquement la définition fondamentale de l'endomorphisme symétrique : $\langle Ax, y \rangle = \langle x, Ay \rangle$ pour tout couple de vecteurs.



\begin{lstlisting}
def dot_product(x, y):
    assert len(x) == len(y), "Les vecteurs doivent être de même dimension."
    return sum(x[i] * y[i] for i in range(len(x)))

def mat_vec_mult(A, x):
    rows = len(A)
    cols = len(A[0])
    assert cols == len(x), "Dimensions incompatibles pour la multiplication."
    result = [0.0] * rows
    for i in range(rows):
        result[i] = sum(A[i][j] * x[j] for j in range(cols))
    return result

# Matrice symétrique
A = [[2.0, -1.0],
     [-1.0, 3.0]]
x = [1.5, 2.0]
y = [-0.5, 3.5]

Ax = mat_vec_mult(A, x)
Ay = mat_vec_mult(A, y)

# <Ax, y>
prod_1 = dot_product(Ax, y)
# <x, Ay>
prod_2 = dot_product(x, Ay)

assert abs(prod_1 - prod_2) < 1e-9, "L'égalité <Ax, y> = <x, Ay> n'est pas vérifiée."
print(f"<Ax, y> = {prod_1}")
print(f"<x, Ay> = {prod_2}")
print("TP 02 validé : La propriété fondamentale de l'opérateur symétrique est validée.")
\end{lstlisting}

\subsection*{TP 03 : Orthogonalité des sous-espaces propres}

\textbf{Objectif :}

Calculer les valeurs propres et vecteurs propres d'une matrice symétrique 2x2, et vérifier algorithmiquement que les vecteurs propres associés à des valeurs propres distinctes sont strictement orthogonaux, conformément au Théorème Spectral.



\begin{lstlisting}
import math

def eigen_2x2_sym(A):
    # Pour A = [[a, b], [b, d]], l'équation caractéristique est X^2 - (a+d)X + (ad-b^2) = 0
    a, b = A[0][0], A[0][1]
    d = A[1][1]

    trace = a + d
    det = a*d - b*b
    delta = trace*trace - 4*det

    lambda1 = (trace + math.sqrt(delta)) / 2
    lambda2 = (trace - math.sqrt(delta)) / 2

    # Calcul des vecteurs propres
    # Si a - lambda1 != 0, le vecteur propre est (-b, a - lambda1)
    v1 = [-b, a - lambda1] if abs(b) > 1e-9 else [1, 0]
    v2 = [-b, a - lambda2] if abs(b) > 1e-9 else [0, 1]

    return (lambda1, v1), (lambda2, v2)

def dot_product(x, y):
    return sum(xi * yi for xi, yi in zip(x, y))

A = [[4.0, 1.0],
     [1.0, 2.0]]

(l1, v1), (l2, v2) = eigen_2x2_sym(A)

ortho_check = dot_product(v1, v2)
assert abs(ortho_check) < 1e-9, "Les vecteurs propres ne sont pas orthogonaux."

print(f"Valeur propre 1 : {l1:.4f}, Vecteur propre 1 : {v1}")
print(f"Valeur propre 2 : {l2:.4f}, Vecteur propre 2 : {v2}")
print(f"Produit scalaire <v1, v2> = {ortho_check}")
print("TP 03 validé : L'orthogonalité des sous-espaces propres est confirmée.")
\end{lstlisting}

\subsection*{TP 04 : Algorithme d'Orthogonalisation de Gram-Schmidt}

\textbf{Objectif :}

Implémenter l'algorithme de Gram-Schmidt pour transformer une base de l'espace en une base orthonormée, ce qui est essentiel pour diagonaliser effectivement une matrice symétrique.



\begin{lstlisting}
import math

def dot_product(x, y):
    return sum(xi * yi for xi, yi in zip(x, y))

def norm(x):
    return math.sqrt(dot_product(x, x))

def scale(v, alpha):
    return [xi * alpha for xi in v]

def sub(v1, v2):
    return [x - y for x, y in zip(v1, v2)]

def gram_schmidt(basis):
    orthonormal_basis = []
    for v in basis:
        u = v
        for q in orthonormal_basis:
            # u = v - <v, q> * q
            proj_coef = dot_product(v, q)
            proj_vec = scale(q, proj_coef)
            u = sub(u, proj_vec)
        u_norm = norm(u)
        assert u_norm > 1e-9, "La famille de vecteurs n'est pas libre."
        orthonormal_basis.append(scale(u, 1.0 / u_norm))
    return orthonormal_basis

basis = [[1.0, 1.0, 0.0], [1.0, 3.0, 1.0], [2.0, -1.0, 1.0]]
onb = gram_schmidt(basis)

# Vérification de l'orthonormalité
for i in range(len(onb)):
    assert abs(norm(onb[i]) - 1.0) < 1e-9, "Vecteur non normé."
    for j in range(i + 1, len(onb)):
        assert abs(dot_product(onb[i], onb[j])) < 1e-9, "Vecteurs non orthogonaux."

print("Base orthonormée générée :")
for vec in onb:
    print([f"{x:.4f}" for x in vec])
print("TP 04 validé : L'algorithme de Gram-Schmidt fonctionne parfaitement.")
\end{lstlisting}

\subsection*{TP 05 : Validation des matrices orthogonales (Isométries)}

\textbf{Objectif :}

Créer une matrice de rotation, prouver algorithmiquement qu'elle est orthogonale ($P^T P = I$), et vérifier qu'elle conserve strictement la norme euclidienne des vecteurs (isométrie).



\begin{lstlisting}
import math

def mat_mult(A, B):
    rows_A = len(A)
    cols_A = len(A[0])
    cols_B = len(B[0])
    result = [[0.0] * cols_B for _ in range(rows_A)]
    for i in range(rows_A):
        for j in range(cols_B):
            result[i][j] = sum(A[i][k] * B[k][j] for k in range(cols_A))
    return result

def transpose(A):
    return [[A[i][j] for i in range(len(A))] for j in range(len(A[0]))]

def is_identity(A):
    for i in range(len(A)):
        for j in range(len(A[0])):
            expected = 1.0 if i == j else 0.0
            if abs(A[i][j] - expected) > 1e-9:
                return False
    return True

# Matrice de rotation d'angle theta
theta = math.pi / 6  # 30 degrés
P = [[math.cos(theta), -math.sin(theta)],
     [math.sin(theta), math.cos(theta)]]

P_T = transpose(P)
product = mat_mult(P_T, P)

assert is_identity(product), "La matrice P n'est pas orthogonale (P^T P != I)."

# Conservation de la norme
def norm(x): return math.sqrt(x[0]**2 + x[1]**2)
def apply(A, x): return [A[0][0]*x[0] + A[0][1]*x[1], A[1][0]*x[0] + A[1][1]*x[1]]

v = [3.0, -4.0] # Norme 5
v_rot = apply(P, v)

assert abs(norm(v) - norm(v_rot)) < 1e-9, "La matrice orthogonale ne conserve pas la norme."

print("TP 05 validé : Matrice orthogonale vérifiée et conservation des isométries confirmée.")
\end{lstlisting}

\end{document}
